\documentclass[11pt,a4paper]{ctexart}

\usepackage[margin=2.3cm]{geometry}
\usepackage{amsmath,amssymb}
\usepackage{booktabs}
\usepackage{enumitem}
\usepackage{xcolor}
\usepackage[hidelinks]{hyperref}

\setlength{\parindent}{0pt}
\setlength{\parskip}{6pt}
\emergencystretch=3em
\tolerance=2000

\title{\bfseries AI 心理评估平台技术报告\\[0.4em]
\large\mdseries\textit{Agentic AI Web Assessment --- Short Technical Report}}
\author{王泽芃（Montgomery）}
\date{2026 年 9 月 13 日}

\begin{document}

\maketitle

\section*{摘要 / Abstract}

\noindent\textbf{中文}：本文是「Agentic AI Web Assessment Challenge」的技术报告。我使用 Claude Code（一种 agentic AI 编码工具）设计、构建并部署了一个基于 Web 的交互式心理评估平台，用于测量大五人格（BFI-10 简版）与对人工智能的态度。平台包含知情同意、交互式问卷、自动计分、交互式结果可视化、匿名数据收集、研究者仪表盘与试测反馈模块，并已部署在 Cloudflare Workers 与 D1 上。报告概述了系统设计、测量构念与选择理由、计分逻辑、技术栈、AI 辅助开发过程、试测评估计划、遇到的问题、AI 的错误与纠正方式，以及后续改进方向。

\medskip

\noindent\textbf{English}: This is the technical report for the ``Agentic AI Web Assessment Challenge''. I used Claude Code, an agentic AI coding tool, to design, build, and deploy a web-based interactive assessment platform measuring the Big Five personality (BFI-10) and attitudes toward artificial intelligence. The platform includes informed consent, an interactive questionnaire, automated scoring, interactive result visualization, anonymous data collection, a researcher dashboard, and a pilot-feedback module, deployed on Cloudflare Workers and D1. The report summarizes the system design, the measured constructs and their rationale, the scoring logic, the technology stack, the AI-assisted development process, the pilot-evaluation plan, problems encountered, mistakes made by the AI and how they were corrected, and directions for future improvement.

\section{构建了什么 / What I Built}

\noindent\textbf{中文}：我构建了一个名为「AI 心理评估平台」的 Web 应用，公开部署于 \url{https://ai-assessment-platform.yangyannan2005.workers.dev}。平台包含六个模块：(1) 介绍与知情同意页，说明评估目的、收集的信息、信息用途以及「仅供教育/研究探索、非临床诊断」的声明；(2) 交互式问卷，包含大五人格简版量表（BFI-10，10 题）与对人工智能的态度量表（10 题），均采用 5 点李克特量表；(3) 自动计分系统，将原始作答标准化为 0--100 分；(4) 交互式结果页，展示总体积极态度得分、人格剖面雷达图、各维度条形图与文字解释；(5) 数据收集与研究者仪表盘，展示参与者数量、分数分布、各维度均值、题目层面分布、克隆巴赫 $\alpha$ 系数与相关矩阵；(6) 试测反馈模块，收集 4 项评分与两条开放性意见。源代码托管于 \url{https://github.com/Montgomery0612/ai-assessment-platform-}。

\medskip

\noindent\textbf{English}: I built a web application called ``AI Psychological Assessment Platform'', publicly deployed at \url{https://ai-assessment-platform.yangyannan2005.workers.dev}. It has six modules: (1) an introduction and informed-consent page stating the purpose, the data collected, how the data will be used, and a ``for educational/research exploration, not a clinical diagnosis'' disclaimer; (2) an interactive questionnaire with the BFI-10 Big Five short scale (10 items) and an AI-attitude scale (10 items), both on a 5-point Likert scale; (3) an automated scoring system that standardizes raw responses to 0--100; (4) an interactive results page showing an overall positive-attitude score, a personality radar chart, per-dimension bar charts, and written interpretation; (5) data collection and a researcher dashboard showing participant count, score distributions, dimension means, item-level distributions, Cronbach's $\alpha$, and correlation matrices; and (6) a pilot-feedback module collecting four ratings and two open-ended comments. Source code is hosted at \url{https://github.com/Montgomery0612/ai-assessment-platform-}.

\section{测量的构念 / Constructs Measured}

\noindent\textbf{中文}：平台测量两类构念。人格维度采用大五人格框架（外向性、宜人性、尽责性、神经质、开放性），通过 BFI-10 简版量表测量，每个维度 2 题。态度维度测量「对人工智能的态度」，包含五个分量表：对 AI 的信任、审慎使用、AI 担忧（反向计分）、使用意愿、AI 社会影响看法；其中第 10 题（价值取向）为描述性题目，不计入总分。

\medskip

\noindent\textbf{English}: The platform measures two kinds of constructs. Personality is measured with the Big Five framework (extraversion, agreeableness, conscientiousness, neuroticism, openness) using the BFI-10 short scale, two items per dimension. Attitude is measured as ``attitude toward artificial intelligence'' across five subscales: trust in AI, cautious use, AI concern (reverse-scored), use intention, and perceived social impact of AI; one item on value orientation is descriptive and excluded from scoring.

\section{为什么选择这些构念 / Why These Constructs}

\noindent\textbf{中文}：选择大五人格，是因为它是心理学中成熟且被广泛验证的人格模型之一，BFI-10 是公认的简短测量工具，适合在试测中快速施测。选择「对 AI 的态度」，是因为 AI 工具在当前学习与工作中日益普及，测量使用者对 AI 的信任、担忧与使用意愿既有现实意义，也便于与人格特质做关联分析（例如开放性与 AI 接受度的关系）。

\medskip

\noindent\textbf{English}: I chose the Big Five because it is one of the most mature and widely validated personality models, and the BFI-10 is an established brief instrument suitable for a fast pilot. I chose ``attitude toward AI'' because AI tools are increasingly pervasive in study and work; measuring trust, concern, and use intention is both timely and convenient for exploring relationships with personality traits (for example, openness and AI acceptance).

\section{计分系统如何工作 / How Scoring Works}

\noindent\textbf{中文}：所有题目采用 1--5 分李克特量表。反向计分的题目先用 $6$ 减去原始分进行翻转。每个人格维度取该维度 2 题（翻转后）的平均分 $\bar{x}$，再通过 $((\bar{x}-1)/4)\times 100$ 线性变换为 0--100 分。AI 态度分量表同理：先对分量表内题目求均值，其中「AI 担忧」分量表反向（分数越高代表担忧越少），再标准化为 0--100。总体积极态度为信任、担忧（反向）、使用意愿、社会影响四个分量表均值的平均，再标准化。审慎使用作为描述性指标单独报告，不计入总体。

\medskip

\noindent\textbf{English}: All items use a 1--5 Likert scale. Negatively worded items are reverse-scored by subtracting the raw value from 6. Each personality dimension is the mean of its two (reverse-scored) items, linearly transformed to 0--100 via $((\mathrm{mean}-1)/4)\times 100$. AI subscales follow the same procedure; the ``concern'' subscale is reversed so that higher scores mean less concern. The overall positive-attitude score is the mean of trust, concern (reversed), use intention, and social impact, then standardized. Cautious use is reported separately as a descriptive indicator and excluded from the overall score.

\section{使用的技术 / Technologies Used}

\noindent\textbf{中文}：前端与全栈使用 Next.js 16（App Router）与 React 19，Node.js 运行时。后端通过 Cloudflare Workers 托管（借助 OpenNext 适配器），数据存储使用 Cloudflare D1（SQLite）。部署与资源管理使用 wrangler CLI。本地验证使用 Node 内置测试与脚本。图表（雷达图与条形图）为手写 SVG，无需额外图表库。

\medskip

\noindent\textbf{English}: The front end and full stack use Next.js 16 (App Router) and React 19 on the Node.js runtime. The backend is hosted on Cloudflare Workers (via the OpenNext adapter), with data stored in Cloudflare D1 (SQLite). Wrangler CLI is used for deployment and resource management. Local verification uses Node's built-in tests and scripts. Charts (radar and bar) are hand-written SVG, requiring no extra charting library.

\section{如何使用 Agentic AI / How I Used Agentic AI}

\noindent\textbf{中文}：整个开发过程使用 Claude Code 完成。我将开放性问题分解为若干子任务：需求与数据模型设计、量表配置、计分逻辑、前端页面、结果可视化、管理员仪表盘、部署与调试。AI 生成了大部分代码，我负责审查、测试与纠正。关键交互包括：设计计分函数（标准分与反向计分）、将本地 JSONL 存储迁移到 Cloudflare D1、用 \texttt{wrangler tail} 定位 D1 建表语句错误、修复 OpenNext 部署配置。具体的关键提示与纠错过程见第 9--10 节。

\medskip

\noindent\textbf{English}: The entire development was done with Claude Code. I decomposed the open-ended problem into subtasks: requirements and data-model design, scale configuration, scoring logic, front-end pages, result visualization, the admin dashboard, and deployment/debugging. The AI generated most of the code, while I reviewed, tested, and corrected it. Key interactions include designing the scoring functions (standardization and reverse scoring), migrating local JSONL storage to Cloudflare D1, using \texttt{wrangler tail} to locate a D1 table-creation error, and fixing the OpenNext deployment configuration. Specific prompts and corrections are detailed in Sections 9--10.

\section{试测评估 / Pilot Evaluation}

\noindent\textbf{中文}：平台已公开部署。为验证系统的端到端正确性，我先用 12 份模拟作答完成了系统级验证，确认提交、计分、结果展示、数据入库、统计聚合与反馈链路均正常。正式试测计划邀请至少 10 名独立被试（同学、朋友等）完成测评并填写反馈。届时将分析：(1) 被试能否顺利完成；(2) 题目是否易于理解；(3) 是否存在可用性问题；(4) 计分是否符合预期；(5) 是否出现异常作答模式；(6) 用户喜欢与不喜欢之处。需要说明：截至本报告撰写时，真实被试数据尚未收集，本节将在收集真实数据后补充具体结论。

\medskip

\noindent\textbf{English}: The platform is publicly deployed. To verify end-to-end correctness, I first ran 12 simulated responses as a system-level check, confirming that submission, scoring, result display, storage, statistical aggregation, and feedback all work. The formal pilot will invite at least 10 independent participants (classmates, friends, etc.) to complete the assessment and give feedback. I will then analyze: (1) whether participants could complete the assessment; (2) whether items were understandable; (3) whether usability problems occurred; (4) whether scoring behaved as expected; (5) whether unusual response patterns appeared; and (6) what users liked or disliked. Note: at the time of writing, real pilot data has not yet been collected; this section will be completed with concrete findings once real participant data is available.

\section{遇到的问题 / Problems Encountered}

\noindent\textbf{中文}：主要问题包括：(1) Cloudflare Workers 没有持久文件系统，原来的 JSONL 本地存储无法使用，需迁移到 D1；(2) Cloudflare 的 workers.dev 域名在中国大陆网络下难以直连；(3) Workers 免费版 CPU 时间有限，运行 Next.js 服务端渲染需要付费方案；(4) OpenNext 对 Windows 的支持有限。这些问题分别通过改用 D1、升级到付费计划、并在报告中说明访问限制等方式处理。

\medskip

\noindent\textbf{English}: The main problems were: (1) Cloudflare Workers has no persistent filesystem, so the original local JSONL storage could not be used and had to be migrated to D1; (2) the workers.dev domain is hard to reach directly from mainland-China networks; (3) the Workers free tier has limited CPU time, so running Next.js server-side rendering requires a paid plan; and (4) OpenNext has limited Windows support. These were handled by moving to D1, upgrading to a paid plan, and documenting the access limitations.

\section{AI 犯的错误 / Mistakes the AI Made}

\noindent\textbf{中文}：AI 在开发中犯了若干错误，由我发现并纠正：(1) OpenNext 配置文件误命名为 \texttt{open-next.config.mjs}，实际要求 \texttt{open-next.config.ts}；(2) \texttt{wrangler.jsonc} 缺少 \texttt{main}、\texttt{assets} 与 \texttt{global\_fetch\_strictly\_public} 等必需字段；(3) 用 D1 的 \texttt{exec()} 一次执行多条分号分隔的建表 SQL，导致 \texttt{SQLITE\_ERROR}，改为逐条 \texttt{prepare().run()}；(4) 使用了 Node 专属 API（\texttt{fs}、\texttt{Buffer} 的 base64url、\texttt{crypto} 的 randomUUID），在 Workers 上不可用，需改为可移植实现；(5) 首次部署时未注册 workers.dev 子域；(6) 在源码中硬编码了管理员回退密码，存在安全隐患。

\medskip

\noindent\textbf{English}: The AI made several mistakes that I identified and corrected: (1) the OpenNext config file was misnamed \texttt{open-next.config.mjs} instead of the required \texttt{open-next.config.ts}; (2) \texttt{wrangler.jsonc} was missing required fields (\texttt{main}, \texttt{assets}, \texttt{global\_fetch\_strictly\_public}); (3) it called D1's \texttt{exec()} with multiple semicolon-separated \texttt{CREATE TABLE} statements, causing a \texttt{SQLITE\_ERROR}, fixed by running one \texttt{prepare().run()} per statement; (4) it used Node-specific APIs (\texttt{fs}, \texttt{Buffer} base64url, \texttt{crypto} randomUUID) that are unavailable on Workers, replaced with portable implementations; (5) it deployed before registering a workers.dev subdomain; and (6) it hard-coded a fallback admin password in source, a security concern.

\section{如何验证与纠正 / How I Verified and Corrected}

\noindent\textbf{中文}：我通过多层验证保证正确性：(1) 语法检查 \texttt{node --check}；(2) 编译检查 \texttt{next build} 成功；(3) 运行既有计分单元测试 \texttt{test/scoring.test.js}；(4) 端到端 API 测试：提交、读取结果、提交反馈、统计聚合，逐一核对返回；(5) 用 \texttt{wrangler tail} 抓取线上日志定位 500 错误（D1 建表语句错误）并修复。每次修复后重新部署并复测。计分逻辑也通过手工核对（全部作答 3 分应得 50 分）确认正确。

\medskip

\noindent\textbf{English}: I verified correctness through multiple layers: (1) \texttt{node --check} syntax checks; (2) a successful \texttt{next build}; (3) the existing scoring unit tests (\texttt{test/scoring.test.js}); (4) end-to-end API tests --- submit, read result, submit feedback, aggregate statistics --- checking each response; and (5) \texttt{wrangler tail} to capture live logs and locate the 500 error (the D1 table-creation bug), then fixing it. After each fix I redeployed and re-tested. Scoring was also checked by hand (an all-3 response should yield 50) to confirm correctness.

\section{如果还有一周 / If I Had Another Week}

\noindent\textbf{中文}：如果还有一周，我会：(1) 收集真实被试数据并完成信度、效度与相关分析；(2) 为大陆用户绑定自定义域名或换用可直连的托管；(3) 用哈希密码与更完善的身份认证替代当前的演示认证；(4) 增加移动端适配与无障碍支持；(5) 丰富结果解释（如维度间比较与个性化建议），并基于试测反馈迭代问卷措辞与交互。

\medskip

\noindent\textbf{English}: With another week I would: (1) collect real participant data and complete reliability, validity, and correlation analyses; (2) bind a custom domain or switch to a host directly reachable for mainland-China users; (3) replace the demo authentication with hashed passwords and a proper auth system; (4) improve mobile responsiveness and accessibility; and (5) enrich result interpretation (for example, cross-dimension comparison and personalized suggestions), iterating on item wording and interaction based on pilot feedback.

\end{document}
